% 依赖项（必须放在主文档导言区，即 \begin{document} 之前）：
% \usepackage{tabularray}
% \UseTblrLibrary{varwidth}
% \usepackage{enumitem}
\begingroup
% \nopagebreak[4]
\enlargethispage{1cm}
\footnotesize
\centering

\DefTblrTemplate{conthead-text}{default}{（续表）}
\DefTblrTemplate{contfoot-text}{default}{续下页}

\begin{longtblr}[
  caption = {国产中文智能体评测基准与本土产学研安全体系},
  label   = {tab:chinese_agent_benchmark_security},
]{
  width   = \linewidth,
  colspec = {
    |Q[wd=0.07\linewidth,l,m]
    |Q[wd=0.14\linewidth,l,m]
    |X[1.00,l,m]|
  },
  cells   = {l,m},
  colsep  = 3pt,
  row{1}  = {font=\bfseries},
  rowhead = 1,
  rowsep  = 5pt,
  measure = vbox,
  stretch = -1,
}
\hline
研究方向 & 企业/机构 + 核心学者 & 自研项目、产学合作、商用安全产品 \\
\hline

\SetCell[r=3]{l,m} 中文本土智能体评测与国产化标准体系
& 清华大学THUNLP/AIR；朱军、唐杰、刘知远、孙茂松
& \begin{itemize}[leftmargin=*,topsep=0pt,partopsep=0pt,itemsep=1pt,parsep=0pt]
\item \textbf{工具调用评测基准：}ToolBench工具调用评测基准，覆盖16000多个真实API中文工具场景~\cite{tsinghuaAIInstitute2026,tsinghuaAIR2026}。
\item \textbf{全栈智能体综合评测：}AgentBench八场景全栈智能体综合评测基准，联合俄亥俄州立大学、加州大学伯克利分校共建。
\end{itemize} \\
\cline{2-3}
& \mbox{北京大学}PKU-Alignment；杨耀东、刘健
& \begin{itemize}[leftmargin=*,topsep=0pt,partopsep=0pt,itemsep=1pt,parsep=0pt]
\item \textbf{安全强化学习训练框架：}Safe RLHF、PKU-SafeRLHF训练框架，BeaverTails中文安全偏好数据集；
\item \textbf{全模态对齐训练框架：}Aligner、ProgressGym等后训练对齐系列成果，开源align-anything全模态对齐训练、数据和测评框架~\cite{pkuAlignAnything2026}。
\end{itemize} \\
\cline{2-3}
& 上海交通大学；谷大武（LoCCS）、张倬胜（AI安全实验室）
& \begin{itemize}[leftmargin=*,topsep=0pt,partopsep=0pt,itemsep=1pt,parsep=0pt]
\item \textbf{密码学与系统安全底座：}谷大武团队围绕密码理论、软硬件与系统安全展开研究，为智能体运行环境提供底层安全基础~\cite{sjtuLoCCS2026}。
\item \textbf{智能体行为安全评测基准：}张倬胜团队提出R-Judge智能体行为安全评测基准，覆盖5类应用场景、27个风险场景与10类风险类型~\cite{sjtuAISecLab2026}。
\end{itemize} \\
\hline

\SetCell[r=3]{l,m} 我国顶尖高校开源智能体对齐安全体系
& 浙江大学区块链与数据安全全国重点实验室；任奎
& \begin{itemize}[leftmargin=*,topsep=0pt,partopsep=0pt,itemsep=1pt,parsep=0pt]
\item \textbf{通用AI安全评测平台：}AIcertAI安全评测平台，覆盖数据质量、算法安全、模型安全、框架安全与系统安全检测，已评测LLaMA、Gemma、Qwen等35个开源大模型，并在淘宝直播风控、中车株洲所等场景形成示范应用~\cite{aicertZJU2026}。
\item \textbf{企业级安全联合实验室：}浙江大学--阿里巴巴集团AI安全联合实验室将智能体安全评测列为首期核心研究方向，支撑企业级大模型Agent红队测试与安全治理~\cite{zjuAlibabaAISafetyLab2025}。
\end{itemize} \\
\cline{2-3}
& 北京大学PKU-Alignment；杨耀东、刘健
& \begin{itemize}[leftmargin=*,topsep=0pt,partopsep=0pt,itemsep=1pt,parsep=0pt]
\item \textbf{后训练安全对齐生态：}维护Safe RLHF、PKU-SafeRLHF、BeaverTails等安全对齐训练与数据体系，构成我国较早的可复现安全强化学习框架~\cite{pkuAlignAnything2026}。
\end{itemize} \\
\cline{2-3}
& 上海交通大学AI安全实验室；张倬胜
& \begin{itemize}[leftmargin=*,topsep=0pt,partopsep=0pt,itemsep=1pt,parsep=0pt]
\item \textbf{自然语言处理与自主智能体安全：}专注自然语言处理、预训练语言模型与自主智能体安全，R-Judge基准支撑智能体行为风险识别研究~\cite{sjtuAISecLab2026}。
\end{itemize} \\
\hline

\SetCell[r=3]{l,m} 我国科研院所、产业实验室行业智能体安全落地
& 中国科学院自动化所；曾毅、王飞跃团队
& \begin{itemize}[leftmargin=*,topsep=0pt,partopsep=0pt,itemsep=1pt,parsep=0pt]
\item \textbf{多智能体博弈攻防评测平台：}PandaGuard（灵御）多智能体博弈攻防评测平台，集成19种攻击算法与12种防御机制~\cite{pandaguardCASIA2026}。
\item \textbf{多智能体决策开放平台：}JidiAI多智能体开源开放平台面向智能体博弈决策领域，吸引我国外50余所高校机构注册使用~\cite{jidiAI2026}。
\end{itemize} \\
\cline{2-3}
& 上海AI实验室；安全可信AI团队
& \begin{itemize}[leftmargin=*,topsep=0pt,partopsep=0pt,itemsep=1pt,parsep=0pt]
\item \textbf{产业级安全智能体操作系统：}Intern-Shannon（书安）产业级安全智能体操作系统，以“安全即服务”模式支撑高安全场景部署~\cite{intershannon2026}。
\item \textbf{安全能力协同演化与科研智能体：}SafeWork-R1通过SafeLadder框架实现安全与智能协同进化，依托InternLM/OpenMMLab开源底座；InternAgent构建面向自主科学发现的闭环多智能体框架~\cite{safework2025,internAgentSHLab2026}。
\end{itemize} \\
\cline{2-3}
& 华为诺亚方舟实验室；郝建业（决策推理实验室），上交陈海波（IPADS实验室）
& \begin{itemize}[leftmargin=*,topsep=0pt,partopsep=0pt,itemsep=1pt,parsep=0pt]
\item \textbf{自动驾驶多智能体仿真平台：}SMARTS大规模多智能体强化学习仿真平台，支持数百至数千个智能体在真实交通场景中的交互学习~\cite{noahArkSMARTS2026}。
\item \textbf{可信执行环境与操作系统级支撑：}sNPU可信执行环境研究被ISCA 2024接收，陈海波担任OpenHarmony技术指导委员会创始主席，IPADS与华为在鸿蒙操作系统上保持长期合作关系~\cite{ipadsSJTU2026}。
\end{itemize} \\
\hline
\end{longtblr}
\endgroup
